\documentclass[11pt]{article}
% Preámbulo compartido para los modelos de parcial ACN
\usepackage[utf8]{inputenc}
\usepackage[T1]{fontenc}
\usepackage[spanish]{babel}
\usepackage{amsmath,amssymb}
\usepackage[margin=2.3cm]{geometry}
\usepackage{enumitem}
\usepackage{fancyhdr}
\usepackage{array}

\pagestyle{fancy}
\fancyhf{}
\lhead{Aplicaciones Computacionales en Negocios --- UTDT}
\rhead{\thepage}
\renewcommand{\headrulewidth}{0.4pt}

% Encabezado tipo examen
\newcommand{\encabezado}[1]{%
  \begin{center}
    {\Large\bfseries Examen parcial}\\[2pt]
    {\large Aplicaciones Computacionales en Negocios --- UTDT}\\[2pt]
    {\normalsize #1}\\[8pt]
  \end{center}
  \noindent\fbox{\parbox{\linewidth}{\centering\bfseries RESPONDER EXCLUSIVAMENTE DEBAJO DEL ESPACIO PROVISTO BAJO CADA PREGUNTA}}
  \vspace{6pt}

  \noindent Nombre y apellido: \underline{\hspace{7cm}} \hfill Legajo: \underline{\hspace{3.5cm}}
  \vspace{6pt}
  \hrule
  \vspace{10pt}
}

% Puntaje en negrita
\newcommand{\pts}[1]{\textbf{[#1 puntos]}\ }

% Espacio de respuesta
\newcommand{\espacio}[1]{\vspace{#1}}


\begin{document}
\encabezado{Modelo 2}

% ------------------------------------------------------------------
\noindent\textbf{1)} Para cada uno de los siguientes procesos estocásticos, respondé: ¿qué rol cumple cada uno de sus parámetros? ¿Cómo es el gráfico de su evolución temporal si $\sigma = 0$?

\medskip
\pts{15} Random walk aritmético con drift:
\[
S_{i+1} = S_i + \mu\,\Delta + \sigma\sqrt{\Delta}\,Z_{i+1}
\]
¿Cómo es su distribución terminal para un horizonte de $N$ pasos ($N$ grande) si $\Delta \ll 1$? Hacé un gráfico (no la fórmula) de esa distribución.

\espacio{4.5cm}

\pts{15} Random walk geométrico (multiplicativo):
\[
S_{i+1} = S_i\left(1 + \mu\,\Delta + \sigma\sqrt{\Delta}\,Z_{i+1}\right)
\]
¿Por qué este modelo es más apropiado que el aritmético para el precio de una acción? ¿Qué forma tiene su distribución terminal?

\espacio{4.5cm}

\pts{10} Proceso con reversión a la media:
\[
S_{i+1} = S_i + \kappa\,(\theta - S_i)\,\Delta + \sigma\sqrt{\Delta}\,Z_{i+1}
\]
¿Por qué decimos que este proceso tiene reversión a la media? ¿Qué sucede para valores de $S_i$ por encima y por debajo de $\theta$? ¿Cuál es la distribución de largo plazo si $\mu = 0$ y $\sigma \neq 0$?

\espacio{4.5cm}

% ------------------------------------------------------------------
\noindent\textbf{2)} Un call center recibe llamados que llegan según un proceso de Poisson con tasa constante $\lambda$ llamados por minuto. Suponé que tenés disponible una función para generar variables uniformes $[0,1]$.

\medskip
\pts{15} Explicá de manera precisa (pseudocódigo y palabras) \emph{dos} algoritmos alternativos para simular la secuencia de llamados durante una hora: (a) uno exacto, sampleando los tiempos entre eventos; (b) uno aproximado, discretizando el tiempo en pasos chicos. ¿Qué condición debe cumplir el tamaño del paso $d$ en el método aproximado y por qué?

\espacio{6.5cm}

\pts{10} Supongamos ahora que la tasa depende de la hora del día, $\lambda(t)$, con $\lambda(t) < \lambda_{\max}$. ¿Cómo modificás el algoritmo aproximado para tener en cuenta esto?

\espacio{4cm}

% ------------------------------------------------------------------
\noindent\textbf{3)} Con el random walk geométrico ya calibrado del punto 1, con parámetros $\mu$ y $\sigma$ diarios y $\Delta$ correspondiente a un día de trading, quiero valuar el riesgo de una acción a un horizonte de un mes ($\sim 21$ días de trading).

\medskip
\pts{20} Describí con precisión (pseudocódigo y palabras) cómo estimarías vía simulación de Monte Carlo: la media $M$ del precio terminal, su desviación estándar $D$, y la probabilidad $P$ de que el precio caiga más de un $10\%$ respecto de hoy. Indicá cómo acompañarías cada estimador con su error estándar.

\espacio{7cm}

\end{document}
